\chap{需求与语义约束：在持续变化中保持软件方向}
\label{sec:intent}
\subsection{意图需要联系正在运行的软件}
人的需求常以不完整的方式进入系统。“取消要可靠”“让我知道现在发生了什么”表达了关切，却没有确定所有状态与边界。智能开发系统需要把这些表达放回具体软件：谁在何时执行任务，哪些操作会重复，什么结果会让人相信任务已经停止。目标解释是逐步形成可行动理解的过程，其中既有可以依据已有信息解决的细节，也有必须由人决定的取舍。

有效澄清应围绕会改变行为的差异。例如，“取消”可能表示不再接收新步骤、尽快停止计算、撤销已产生的副作用，或仅关闭界面中的等待。这些含义对资源、数据和用户预期有不同影响。Agent 可以先检查已有契约和使用场景，再提出足够具体的选项。要求人一次写出完整规范，容易把工程分析工作重新推回给人。

意图还包括长期偏好。有的项目优先缩短故障恢复时间，有的项目对数据一致性有严格要求，同一项目不同功能也可能采用不同取舍。系统应保存偏好的范围和理由。例如，“局部修复通过定向验证即可交付”适用于已授权的维护范围，不能据此推出所有数据迁移都可以立即执行。明确范围让团队能够自主行动，并减少重复确认。

\subsection{语义规范的内容与颗粒度}
语义规范用人和 Agent 都能理解的语言说明目标、行为约束与成功含义。它可以包含实例、状态定义、接口契约、异常情境及重要的非功能要求。自然语言的优势是能解释动机和取舍；其含糊之处则需要通过具体情境和检查方式收敛。过于抽象的愿望难以指导实现，过于细碎的实现指令又容易随代码变化而失效。

对于取消功能，一份有用的规范可以说明：取消请求被接收后，用户立即看到处理中状态；只有在执行器确认停止或确定任务已经完成之后，才显示实际终态；重复请求应得到一致的解释；可保留的部分结果仍能访问。这些要求表达用户可依赖的行为，而具体由轮询、事件订阅或其他机制实现，可以由模块负责人决定。

\begin{table}[htbp]\centering\small
\caption{从抽象需求到可讨论的语义约束}\label{tab:intent}
\begin{tabularx}{\linewidth}{p{2.7cm}YY}\toprule
初始表达 & 需要明确的含义 & 可用于核对的情境\\\midrule
取消要可靠 & 请求接收与停止完成的区别 & 执行器延迟响应时界面如何变化\\
不要丢失结果 & 哪些中间产物可保留、何时可见 & 部分结果已写出时取消\\
操作要简单 & 用户需要几步完成并理解结果 & 不熟悉内部机制的用户首次使用\\
修复要及时 & 允许的交付范围与必要验证 & 旧版本异常时如何送达局部修复\\\bottomrule
\end{tabularx}\end{table}

规范的颗粒度可以沿职责组织。用户任务说明端到端成功含义，模块契约规定局部责任，稳定规则由代码检查反复核对。不同粒度的材料需要有明确关系，否则同一个词可能在多个文件中逐渐产生不同含义。维护规范的目标是帮助下一次判断，而非使文档条目持续增长。

\subsection{让目标、实现和证据可以互相定位}
每次改进应能够回答三个问题：它服务于什么目标，改变了哪些对象，依据什么判断已经成立。对象可以是模块、接口、用户任务或界面状态；证据可以是代码测试、真实使用记录、运行观察或人类反馈。稳定的关联使成员能够按需深入，而不必重复扫描全部资料。

关联也需要反向可用。一个测试失败时，应能找到相关契约与责任模块；一个用户指出界面含义不清时，应能定位它显示的实际状态；某个协议改变时，应能识别哪些使用任务和投影受到影响。完整关联不要求为每行代码建立手工索引，优先维护那些会影响协作和判断的边界即可。

\fig{03-semantics}{语义约束与工作证据的联系。需求经场景解释形成契约，契约进入实现与验证；出现冲突时，根据目标和证据定位需要修订的部分。}{semantics}

这套联系也为汇报生成提供事实基础。Agent 可以从目标出发选择已经交付的变化，从具体对象获取可展示实例，再关联尚未解决的问题。若汇报只依据成员口头声明，就可能把准备工作当成结果；若只依据代码差异，又可能忽略用户体验。需要将内容、运行和使用证据组合起来解释进展。

\subsection{一个可使用的行为说明实例}
为了让语义材料能够直接支持协作，可以围绕一个行为对象组织简短说明。对象名为“运行中任务的取消”，目标是让用户可靠表达停止意图并理解实际结果；参与方包括用户界面、任务服务、执行器和结果存储；前置条件是任务已经创建且用户具备相应操作权限。行为说明随后列出接收、等待、停止确认及自然完成等情境。

这份说明还应包含可观察条件。请求接收可以通过持久记录核对，停止确认可以通过执行器状态与清理条件核对，用户理解则需要通过实际任务观察。不同证据支持不同结论，不能仅以接口返回成功证明任务已停止。明确这一点后，实现者与验证者可以独立选择合适工具。

当出现“任务取消后保留日志”的新需求时，说明只需更新结果保留部分及其关联情境。取消请求幂等性仍然是原有约束，未变化的部分继续有效。这种按对象与条件组织的材料有助于渐进维护，也为界面生成提供稳定语义。它不要求预先建立复杂的规范语言，最初可以用带明确字段和实例的文本实现。

\subsection{规范、测试与实现发生冲突时}
快速迭代中，规范可能落后于需求，测试可能保留过时断言，实现也可能偏离原意。发现不一致后，不能仅以“最新”判定正确。首先要确认各材料所描述的对象和适用条件是否相同，再检查哪项假设发生了变化。如果用户已改变输出保留要求，那么原先“取消后删除全部结果”的测试需要更新；如果用户意图未变，则当前实现可能存在缺陷。

修订应保留足以解释后续变化的理由。例如，把“已取消”拆成“取消请求已接收”与“停止已确认”，应说明它解决的是用户对实际执行状态的误解。理由可以关联一个复现情境和对应检查。这样，未来成员看到两个状态时，能够理解它们为何存在，也不容易为了界面简洁再次合并它们。

规范维护本身可以委派，但关键目标的变化仍需来自有效的人类输入或既定约束。Agent 可以提出新的解释和方案，并把影响具体展示出来；它不应通过不断调整成功定义，使所有实现都显得正确。评价系统时，也要防止执行者同时任意改写评价标准而掩盖失败。

\subsection{工作状态与知识的不同生命周期}
任务排队、运行日志、当前错误和一次测试结果属于工作记录；模块职责、稳定契约和经过核对的经验具有更长使用周期。将原始运行记录全部归档为长期知识，会增加检索噪声。相反，只保存简短结论而丢弃依据，也会让后来者无法判断结论是否仍适用。

一种可行做法是保留原始记录的入口，并提炼带条件的知识。经验条目可以包含“什么问题下有效”“采取什么处理”“怎样核对”“什么情况下需要重新检查”。同一经验在多个任务中被验证后，其适用范围可能扩大；遇到反例时则需要缩小或修订。知识的价值来自后续使用效果，而非被写入了多少文件。

WALL 类界面适合保存和关联设计点，但面向人的默认视图应根据当前问题组织信息。需要了解架构时，设计点附着在模块和协议上；需要验收功能时，它们附着在行为与结果上。这样，完整信息仍可查询，人的日常参与则围绕有意义的情境展开。

\subsection{把澄清作为持续协作的一部分}
需求通常在看到具体成果后变得更清楚。因此，早期阶段应尽快形成能够被理解和体验的中间成果。它可以是协议示例、状态模拟或简化界面。人对这些对象的反馈，比孤立地回答抽象问题更容易包含准确约束。

持续澄清也意味着并非所有未知都需要立即解决。对当前实现无影响的远期问题可以明确保留；会决定模块边界或数据含义的问题则应优先处理。Agent 需要判断问题的重要性和时机，避免通过大量提问制造新的控制负担。其成功标准是理解逐步变准确，并能支持有质量的行动。
